\chapter*{Introduction}
\label{toc:introduction}

This document describes the \FramaC/\Region plug-in that provides a
modular aliasing and memory access analysis. It deserves several purposes. For
\FramaC users, it can help understand the memory layout by providing a visual
representation of memory accesses through the code. These memory maps can be
visualized from the \FramaC Graphical User Interface. On the other hand, for
\FramaC plug-in developers, it provides a rich collection of aliasing and memory
access properties that can be used to improve the precision of other analyses.
Typically, the \FramaC/\Wp plug-in has a dedicated memory model
(\verb$-wp-model region$) that builds upon the memory maps computed by the
\FramaC/\Region plug-in.

Memory maps computed by \FramaC/\Region are made of \emph{regions} that
denotes disjoint portions of the memory with consistent layout, access kinds and
properties (validity, initialization, etc.). The plug-in computes one memory map
for every function, hence it is a modular analysis. The memory context for every
function is inferred from the code and can be further refined by the user thanks
to dedicated \acsl annotations. The analysis is \emph{optimistic} in the sense
that by default, the analysis assumes that different pointers will point to
disjoint memory regions unless they are explicitly aliased by the code. However,
this might not be consistent with actual caller contexts, for which a
\emph{modular} analysis has no insight. To adapt to those (hopefully rare)
situations, the user shall provide dedicated region annotations to let the
\FramaC/\Region plug-in being less \emph{optimistic} when needed.

Hence, for the memory maps computed by \FramaC/\Region to be consistent,
there are some side conditions to be verified. The memory region can generate
them as additional \acsl annotations (\verb$-region$) that will instrument the
code in such a way that, when satisfied, the memory maps and their associated
properties are actually valid. Those annotations play the same role as the ones
generated by the \FramaC/\rte plug-in, although they are more precise.
Alternatively, other plug-ins can use the \FramaC/\Region API to obtain
the side-conditions to be validated.

The document is structured as follows: Chapter~\ref{toc:userguide} provides a
user guide to the \Region plug-in via its command-line interface and its
graphical interface in the GUI. Chapter~\ref{toc:reference} provides
a reference guide on the generated side-conditions.
